\documentclass[conference]{IEEEtran}
\usepackage[utf8]{inputenc}
\usepackage[T1]{fontenc}
\usepackage{cite}
\usepackage{amsmath,amssymb}
\usepackage{graphicx}
\usepackage{booktabs}
\usepackage{url}

\title{Verifier‑Assisted versus One‑Shot LLM Intent\ensuremath{\rightarrow}Configuration: A Paired NetOpsTraceBench Study}

\author{
\IEEEauthorblockN{Autonomous AI Researcher}
\IEEEauthorblockA{CNSM 2026 Agentic AI Researcher Track}
}

\begin{document}

\maketitle

\begin{abstract}
We preregistered and executed NetOpsTraceBench, a paired comparison of a one‑shot LLM intent\ensuremath{\rightarrow}configuration pipeline and a verifier‑assisted generate--validate--repair pipeline on 1,200 intent\ensuremath{\rightarrow}config episodes. Using gpt-5-mini, a deterministic NetOps validator, and Mininet‑style emulation, we applied a preregistered binary episode success rule (validator accepts final config AND emulator shows no availability degradation above tolerance AND no automatic rollback). Of 1,200 planned paired tasks we completed 1,199 pairs. Baseline success = 1,196/1,199 (0.9974979149); guarded success = 1,199/1,199 (1.0000000000). The paired‑success‑rate difference (guarded - baseline) = 0.0025020851 (95\% bootstrap percentile CI [0.0000, 0.0058381985], bootstrap\_seed=1729, resamples=10,000); exact McNemar two‑sided p = 0.25. The preregistered power target sought a 5 percentage‑point minimum detectable effect; given the observed near‑ceiling baseline, the study lacked power to detect much smaller absolute gains. Representative validator traces and provider audit logs show repair was invoked only three times. Under these controlled emulation and task‑corpus conditions we find no statistically significant advantage for the verifier‑assisted pipeline; we report artifacts, analytic deviation (Wilcoxon \ensuremath{\rightarrow} exact McNemar for binary outcomes), and detailed execution accounting to support reproducible interpretation.
\end{abstract}

\section{Introduction}
Generative large language models (LLMs) are increasingly proposed to translate operator intent expressed in natural language into executable device configurations for network operations (NetOps). Systems and survey literature emphasize that operational reliability depends critically on orchestration, deterministic validators, and assurance infrastructure rather than on an LLM alone [1,2]. Generate--validate--repair (verifier‑assisted) architectures are a canonical mitigation strategy: an LLM proposes candidate outputs, an automated verifier checks constraints and safety, and an automated repair loop iterates until the candidate satisfies verifier rules [3--5]. Despite conceptual support, publicly archived, preregistered, paired evidence measuring whether verifier assistance yields an operationally meaningful improvement (rollback‑aware binary success) relative to one‑shot generation is sparse.

We address this gap with NetOpsTraceBench: a preregistered, paired, emulator‑grounded study comparing (A) baseline one‑shot NL\ensuremath{\rightarrow}config generation and (B) a guarded generate--validate--repair pipeline using the same shared initial candidate. The study asked: Can a reproducible, workflow‑centered benchmark of paired intent\ensuremath{\rightarrow}config episodes produce an objective paired binary success metric that reliably distinguishes verifier‑assisted pipelines from one‑shot generation across two simulated topologies? Our contributions are (1) a preregistered experimental design and full execution artifacts (study ID rX4f7-1); (2) a paired analysis on 1,199 completed pairs using gpt-5-mini and a deterministic validator; (3) transparent reporting of an analytic deviation (final confirmatory test used an exact McNemar test appropriate for paired binary outcomes) and execution accounting that explains why the guarded pipeline rarely required repair.

\section{Related Work}
Recent survey and system literature frames LLMs as components inside broader AIOps/NetOps workflows and repeatedly emphasizes that correctness and trustworthiness are properties of the whole orchestration stack (model + verifier + provenance + sandboxed execution), not of the base LLM alone [1--3]. Multiple community reviews urge workflow‑centered evaluation---trace quality, validator interactions, and rollback‑aware metrics---rather than isolated one‑shot QA tests; these reviews motivate reproducible, audit‑ready benchmarks that combine generation, deterministic validation, and emulation [1,2,4].

Concrete verifier‑assisted approaches appear across adjacent domains and NetOps prototypes. Work on NL\ensuremath{\rightarrow}formal translations and CTL/specification generation demonstrates generator+verifier loops where formal checks drive iterative repairs; similarly, topology‑compilation systems translate high‑level security intent into executable topologies with compliance checks and constraint enforcement [5,6]. Prototype efforts in the configuration space (text\ensuremath{\rightarrow}config generators and misconfiguration detectors) and domain‑specific "LLM firewalls" illustrate practical patterns: an LLM proposes a candidate, an automated validator enforces syntactic and semantic constraints, and an orchestrator optionally triggers constrained repair or rejects the proposal [7--10]. These systems commonly instrument deterministic parsers and domain rules to produce binary verifier outcomes (accept/reject) and produce their own evaluation sets for demonstration; as a result, reported gains are often conditional on task selection and validator strictness.

Empirical and conceptual studies document LLM failure modes that motivate guarded architectures. Recent analyses identify hallucination, unit/semantic errors, overconfidence/miscalibration, provenance loss, and susceptibility to prompt‑injection as recurrent risks when LLMs propose executable artifacts---risks that an external verifier or provenance linker can detect or mitigate in some instances but not eliminate universally [11--14]. Several papers demonstrate that verifier‑guided repair reduces syntactic and some semantic errors in constrained settings, but they also highlight sensitivity to the verifier's rule set (strictness) and to the representativeness of evaluation examples. In short, prior work shows feasibility of generate--validate‑repair patterns but also underscores the dependence of measured improvement on corpus difficulty, adversarial content, and verifier configuration.

Related benchmarking and dataset discussions emphasize gaps that NetOpsTraceBench targets. System papers often construct private or narrowly scoped held‑out sets calibrated to their prototype; TopoIntent and other recent tools note the absence of large public, vendor‑diverse intent\ensuremath{\rightarrow}config datasets with execution ground truth and rollback semantics, which complicates cross‑system comparisons [6,7,9]. Surveys therefore call for standardized, workflow‑centered benchmarks that (a) record full provider and validator traces, (b) embed emulator execution to check operational side effects (availability/rollback), and (c) preregister analysis choices to reduce researcher degrees of freedom [1,2]. NetOpsTraceBench operationalizes these recommendations by preregistering inclusion/exclusion rules, seeds, a binary, rollback‑aware episode success metric, and by archiving raw model outputs, validator traces, and execution manifests so that the effect of verifier design and task corpus can be examined reproducibly.

Taken together, the literature supports the conceptual value of verifier scaffolding while also documenting two recurring caveats that our study made explicit in design and interpretation: (i) measured gains depend strongly on task difficulty and verifier thresholds, and (ii) many demonstrations use constructed datasets and domain‑specific validators, which limits generality. NetOpsTraceBench contributes a preregistered, artifact‑rich paired evaluation that makes these dependencies observable (repair invocation counts, contingency discordances, and emulator rollback checks) rather than inferred, enabling follow‑up sensitivity studies that vary corpus adversariality and verifier strictness.

\section{Methodology}
Preregistration and scope: The study preregistration (rX4f7-1; preregistration SHA‑256 = 77abe8b4\allowbreak{}93bde444\allowbreak{}84e0b3d0\allowbreak{}2cfd7f9e\allowbreak{}c1621ff7\allowbreak{}e7847889\allowbreak{}082ae177\allowbreak{}a7ed5712) froze task selection, inclusion/exclusion rules, contamination thresholds, stopping rules, seeds, analysis plan, and failure treatments before execution. NetOpsTraceBench\_v1 comprises 1,200 paired intent\ensuremath{\rightarrow}config tasks split into two topology clusters (edge = 600, core = 600). Each task is public synthetic or consented with deterministic provenance recorded. Tasks are vendor‑agnostic canonical sequences rendered into device configurations for emulator execution. Inclusion/exclusion rules were deterministic: pairs were excluded if contamination checks flagged training‑data overlap above thresholds or if both episodes failed to produce parsable vendor renderings after one automated re‑execution attempt; exclusions are logged and reported.

Pipelines and instrumentation: We compared (A) Baseline one‑shot --- a single LLM call producing an initial candidate configuration; (B) Guarded verifier‑assisted --- the same initial candidate plus at most one automated repair iteration triggered when the deterministic validator rejected the candidate. Execution used requested\_model = gpt-5-mini (provider recorded as openai\_responses) orchestrated by adapter hosted\_netops\_gvr\_v1. Deterministic\_netops\_validator\_v5 handled parsing, intent‑constraint checks, and emulator preflight; the task generator was deterministic\_netops\_task\_generator\_v5. The execution manifest records model\_calls\_used = 1,203 and maximum\_model\_calls = 2,400.

Success metric and confirmatory analysis: The preregistered binary episode success rule required (1) validator acceptance of the final configuration, (2) emulator execution with no availability degradation above the preregistered tolerance, and (3) no automatic rollback during simulated execution. The preregistration specified a nonparametric paired procedure (Wilcoxon signed‑rank) for paired outcomes; because realized outcomes are binary, the archived confirmatory analysis used an exact McNemar two‑sided test on the paired 2\ensuremath{\times}2 contingency table and computed a paired bootstrap percentile 95\% CI for the paired success‑rate difference using frozen bootstrap\_seed = 1729 and 10,000 resamples. This post‑preregistration analytic choice is documented in archived artifacts (analysis/paired\_contingency\_table.csv, analysis/results.json, analysis/analysis\_log.jsonl). The primary estimand is paired\_success\_rate\_difference\_guarded\_minus\_baseline. The preregistered power plan targeted a minimum detectable effect (MDE) of 5 percentage points (absolute) at two‑sided \ensuremath{\alpha}=0.05 and power 0.8 given assumed baseline \ensuremath{\approx}0.70; the power rationale and sensitivity scenarios are recorded in the preregistration.

\section{Results}
Execution and primary accounting (archived artifacts). The run attempted the preregistered 2,400 episodes (1,200 paired tasks) and completed 2,398 episodes (completed\_pairs = 1,199; failed\_episode\_count = 2). Provider audit fields recorded hosted\_model\_call\_event\_count = 1,203, completed\_hosted\_model\_call\_event\_count = 1,202, and failed\_hosted\_model\_call\_event\_count = 1. Stage-level accounting shows shared\_initial\_generation = 1,200 and repair = 3 (execution/provider audit and deterministic\_reconciliation.json). Artifact\_count for the full run = 8,408 and the run completed at 2026-08-30T03:15:12.141890+00:00 (execution manifest).

Primary confirmatory outcomes (archived analysis outputs). Using the preregistered binary success rule, the archived condition summary (analysis/condition\_summary.csv) reports baseline\_success\_count = 1,196 of 1,199 complete pairs (baseline\_success\_rate = 0.9974979149291076) and guarded\_success\_count = 1,199 of 1,199 (guarded\_success\_rate = 1.0). The paired contingency table (analysis/paired\_contingency\_table.csv) is n11 = 1,196 (both succeeded), n10 = 3 (guarded succeeded, baseline failed), n01 = 0 (baseline succeeded, guarded failed), n00 = 0 (both failed). The point estimate for the paired success‑rate difference (guarded - baseline) = 0.002502085070892446. The archived confirmatory exact McNemar two‑sided p-value = 0.25. The paired bootstrap percentile 95\% confidence interval computed with bootstrap\_seed = 1729 and 10,000 resamples is [0.0000, 0.005838198498748957] (analysis/results.json and analysis/analysis\_log.jsonl).

Missingness and failed-call accounting. The preregistered missingness summary and deterministic reconciliation artifacts report complete\_pair\_count = 1,199, failed\_baseline\_episodes = 1, failed\_guarded\_episodes = 1, and both\_missing\_pairs = 1 (analysis/missingness\_summary.csv; deterministic\_reconciliation.json). No pairs were excluded for contamination under the preregistered thresholds (analysis/contamination\_summary.csv empty). The run therefore uses the complete\_pair\_primary treatment as planned.

Repair invocation and provider audit diagnostics. Repair-stage calls were invoked only three times (stage\_counts.repair = 3 in provider\_call\_audit). Provider audit also shows no cross-condition cache reuse (cross\_condition\_cache\_reuse\_observed = false) and confirms the intended shared\_initial\_generation semantics (shared initial candidate used for paired episodes). The low repair count is corroborated by the scoring artifacts: representative scoring JSONs (e.g., execution/scoring/task-000001-baseline.json and execution/scoring/task-000001-guarded.json) show validator\_trace.validation\_after.valid = true and validator\_trace.repair\_applied = false for typical episodes. In other words, the deterministic validator accepted the initial shared candidate in the vast majority of tasks, leaving only three guarded episodes that required a repair call.

Representative per-task diagnostics (artifact-grounded). Representative task\_manifest entries and scoring traces illustrate typical episode structure and validator acceptance criteria. For example, pair-000001 (execution/scoring/task-000001-baseline.json) had parsed\_command\_count = 4, required\_command\_count = 4, observed\_touched\_field\_count = 3, and validation\_after.normalized\_configuration matching the reference; validation\_after.valid = true and repair\_applied = false. A harder task example (pair-000002) shows difficulty.level = "hard" with changed\_interface\_count = 2 and required\_command\_count = 2, yet its initial candidate was still validated as valid by the deterministic validator. These per-episode fields (parsed\_command\_count, required\_command\_count, difficulty pattern) are recorded in the representative scoring artifacts and demonstrate that many tasks contained explicit required sequences which the generator reproduced correctly.

Interpretation of discordant pairs and CI. The observed discordant pattern (n10 = 3, n01 = 0) implies that all discordant evidence favored the guarded pipeline; however, the absolute number of discordant pairs is tiny relative to the total sample. The bootstrap CI lower bound is exactly 0.0000 (percentile method) and the upper bound \ensuremath{\approx} 0.00584, indicating that the data are compatible with a small positive absolute gain but not with the preregistered target MDE (5 percentage points). The exact McNemar p = 0.25 formalizes that, under these discordant counts, the test does not reject the null of no paired difference at \ensuremath{\alpha} = 0.05.

Why repair was rarely exercised. Two archived facts explain the infrequent repair calls: (1) the synthetic task generator produced tasks with explicit required sequences and constrained allowed field changes (deterministic\_netops\_task\_generator\_v5 task\_payloads record required\_sequence entries), increasing the chance that a straightforward canonical candidate satisfies intent constraints; (2) the deterministic validator configuration (deterministic\_netops\_validator\_v5) accepted initial candidates when all intent constraints and preflight checks matched the rendered candidate. Provider audit stage\_counts (shared\_initial\_generation = 1,200, repair = 3) and scoring artifacts jointly document these operational mechanisms and support the conclusion that the guarded pipeline had little opportunity to improve over an already‑high baseline under the chosen corpus and validator strictness.

Ancillary archived metrics. The analysis artifact set includes analysis/paired\_difference\_figure.svg (visual CI and point estimate), analysis/paired\_contingency\_table.csv (full contingency counts), and analysis/analysis\_log.jsonl (bootstrap parameters: resamples = 10,000; seed = 1729). The execution manifest archives timing, provider call counts, and the preregistration SHA‑256 reference; these artifacts are sufficient to reproduce the confirmatory estimates and the diagnostic counts reported above.

\section{Discussion}
Statistical and operational interpretation: The exact McNemar test and bootstrap CI both indicate no statistically significant advantage for the guarded pipeline under this corpus and validator configuration (p = 0.25; 95\% CI includes zero). Two operational explanations are supported by archived artifacts. First, near‑ceiling baseline performance (\ensuremath{\approx}99.75\%) constrained absolute headroom for improvement; the preregistered MDE was 5 percentage points, so the run was not powered to detect the observed \ensuremath{\approx}0.25 percentage‑point effect. The archived CI and discordant counts (n10 = 3, n01 = 0) quantify this limitation and show that detection of much smaller practical gains would require substantially more discordant pairs or a task corpus with lower baseline success.

Second, corpus and verifier design materially affected repair invocation. Representative task\_manifest entries and scoring traces show tasks were synthetic with explicit required sequences and conservative initial generators (deterministic\_netops\_task\_generator\_v5), and the deterministic validator (deterministic\_netops\_validator\_v5) accepted the majority of initial candidates. Those design choices---synthetic tasks yielding clear required sequences and a validator configuration that enforces the preregistered intent constraints---likely reduced baseline difficulty and therefore the frequency of repairs. The provider audit (repair = 3) and scoring artifacts (validation\_after.valid = true for typical episodes) jointly support this interpretation.

Threats to validity and externality: Internal validity is strong for deterministic components (frozen seeds, deterministic task generation, and automated exclusion rules), but construct validity depends on whether the binary success rule (validator acceptance + emulator availability tolerance + no rollback) maps to operationally meaningful correctness across production environments. External validity is limited: only one requested model family (gpt-5-mini) was exercised, task examples were synthetic/consented and vendor‑agnostic, and emulation (Mininet‑style) may not capture vendor idiosyncrasies and routing convergence phenomena. Conclusion validity: the very low discordant counts imply limited power to detect small absolute differences; the preregistered MDE (5 pp) should be considered when interpreting the null result.

Operational implications: When initial generation already meets deterministic verifier criteria for common, well‑specified tasks, generate--validate--repair will be infrequently exercised and measured gains on a binary success metric will be small. To evaluate verifier value, future benchmarks should increase task ambiguity/adversarial content, vary verifier strictness, and exercise multiple model families. The archived artifacts and traces provide a reproducible starting point to perform such sensitivity studies.

\section{Limitations}
\begin{itemize}
\item Single model family tested (gpt-5-mini); results do not imply generality across other LLMs or fine‑tuned variants.
\item Task corpus skew: synthetic/consented tasks yielded near‑ceiling baseline success (\ensuremath{\approx}99.75\%), limiting headroom and statistical power to detect practical improvements by guarded pipelines.
\item Emulator fidelity: Mininet‑style emulation and deterministic validators approximate production behavior but may not capture all deployment failure modes (routing convergence, vendor idiosyncrasies, transient telemetry).
\item Verifier configuration dependence: deterministic\_netops\_validator\_v5 acceptance thresholds and parsing rules materially affect outcomes; different verifier strictness could change repair invocation rates and measured gains.
\item Analytic deviation documented: the preregistration specified a Wilcoxon signed‑rank paired procedure; the archived confirmatory analysis used an exact McNemar test because outcomes were binary. This deviation is recorded in the archived analysis artifacts and in the Disclosure Statement above.
\item Operational external validity: absence of heterogeneous vendor renderers, real production templates, and live rollouts constrain direct production generalization.
\end{itemize}

\section{Conclusion}
We preregistered and executed NetOpsTraceBench, a paired, emulator‑grounded study comparing one‑shot and verifier‑assisted NL\ensuremath{\rightarrow}configuration pipelines on 1,199 completed pairs. Under our synthetic task corpus, deterministic validator settings, and gpt-5-mini as requested, baseline one‑shot generation already achieved near‑ceiling success (\ensuremath{\approx}99.75\%). The guarded pipeline increased absolute success by \ensuremath{\approx}0.25 percentage points (paired difference = 0.0025020851) with exact McNemar two‑sided p = 0.25 and a 95\% bootstrap CI that includes zero. Archived execution manifests, provider audits (model\_calls\_used = 1,203; repairs invoked = 3), representative validator traces, and analysis artifacts support the interpretation that the observed null result reflects limited headroom given task and verifier choices rather than definitive evidence that verifiers are never useful. Future work should intentionally design more challenging and adversarial task sets, vary verifier strictness, and test multiple LLM families to determine conditions under which verifier scaffolding yields operationally meaningful improvements. All run artifacts, analysis code, seeds, and logs are archived to enable exact reproduction and follow‑up sensitivity studies (see Disclosure Statement).

\section*{Disclosure Statement}
Mandatory disclosure (preregistered run rX4f7-1). LLM(s) and provider: requested model = gpt-5-mini via provider recorded as openai\_responses; model configuration archived at execution/model\_configuration.json (requested\_model=gpt-5-mini, max\_output\_tokens=1500). Orchestration/adapter: hosted\_netops\_gvr\_v1 executed the paired benchmark. Deterministic tooling: deterministic\_netops\_validator\_v5 performed canonical parsing, intent‑constraint validation, and emulator preflight checks; deterministic\_netops\_task\_generator\_v5 produced synthetic tasks. Exact initial master prompt: archived at provenance/master\_prompt.txt (immutable reference SHA‑256 = 1872df1e\allowbreak{}1805d2d9\allowbreak{}6940456c\allowbreak{}a016bd66\allowbreak{}5d1d5196\allowbreak{}add77f5a\allowbreak{}cdf1582b\allowbreak{}b39b15ba). Topic constraint and autonomy boundary: the human Research Director selected the topic family and approved the master prompt pre‑lock; after the autonomy lock the AI autonomously formulated the research question, conducted literature review, generated the preregistration, executed the experiment, performed analysis, and wrote and revised the manuscript. Preregistration: NetOpsTraceBench (ID rX4f7-1) preregistration manifest SHA‑256 = 77abe8b4\allowbreak{}93bde444\allowbreak{}84e0b3d0\allowbreak{}2cfd7f9e\allowbreak{}c1621ff7\allowbreak{}e7847889\allowbreak{}082ae177\allowbreak{}a7ed5712; preregistered design (confirmatory hypotheses, inclusion/exclusion, seeds, stopping rules, analysis plan) was frozen before execution. Execution summary: started 2026-08-30T01:21:55.415955+00:00, completed 2026-08-30T03:15:12.141890+00:00; planned episodes 2,400; completed episodes 2,398; completed pairs 1,199; model\_calls\_used = 1,203; repair-stage calls observed = 3. Analysis artifacts and deviation record: the preregistration specified a Wilcoxon signed‑rank paired procedure; the archived executed confirmatory analysis used an exact McNemar two‑sided test on paired binary outcomes plus a paired bootstrap percentile CI. This post‑preregistration analytic choice is documented in the archived analysis artifacts (analysis/paired\_contingency\_table.csv, analysis/results.json, analysis/analysis\_log.jsonl, analysis/paired\_difference\_figure.svg) produced as part of the run that completed at 2026-08-30T03:15:12.141890+00:00. Artifacts and reproducibility: raw model outputs, provider traces, validator traces, task\_manifest, transformation\_manifest, execution logs, analysis code, and all seeds are archived; artifact\_count = 8,408 and full hash manifest path = execution/execution\_manifest.json. Human involvement: pre‑lock collaboration and infrastructure development occurred; no human manual editing, selective revision, or ad hoc text insertion occurred on the manuscript content after the autonomy lock---the AI performed internal autonomous revision cycles only. Technical restarts and deviations: execution manifest records two failed episodes and the analytic deviation described above; all deviations are logged in execution/execution\_log.jsonl. The disclosure above is complete relative to the archived artifacts supplied with this run.

\begin{thebibliography}{99}
\bibitem{ref1} http://hdl.handle.net/20.500.12708/229494
\bibitem{ref2} https://arxiv.org/abs/2608.13389
\bibitem{ref3} http://arxiv.org/abs/2403.09369
\bibitem{ref4} https://doi.org/10.2139/ssrn.6947162
\bibitem{ref5} https://doi.org/10.2139/ssrn.7222278
\bibitem{ref6} https://doi.org/10.20935/acadai8260
\bibitem{ref7} https://doi.org/10.3390/info17030297
\bibitem{ref8} https://doi.org/10.3390/app16010085
\end{thebibliography}

\end{document}
